\begin{abstract}
A central design challenge for future generations of wireless networks is to meet users' ever-increasing demand for capacity, throughput, and connectivity. Recent advances in wireless network design, including ongoing NextG efforts, are driving the transition from Orthogonal Multiple Access (OMA) to Non-Orthogonal Multiple Access (NOMA). This transition aims to support a high number of users near base stations and maximize spectral efficiency. However, NOMA requires sophisticated signal processing to successfully disentangle mutually interfering streams, a process that becomes less efficient as the number of concurrent streams increases. Theoretically optimal Maximum Likelihood (ML) processing could significantly improve performance over traditional methods but soon becomes infeasible due to its computational complexity and processing-time limitations. Thus, the computational capacity of base stations is becoming one of the key limiting factors on performance gains in wireless networks.

Quantum computing is a potential tool to address this computational challenge. It exploits information-processing capabilities based on quantum mechanics to perform calculations that can be intractable by traditional digital methods.

This work presents the application of Quantum Annealing (QA) to enhance ML processing in NOMA systems within a Centralized Radio Access Network (C-RAN) environment. The proposed QA-aided ML decoder serves as a promising alternative to conventional Successive Interference Cancellation (SIC). The analysis incorporates channel randomness and transmit power levels for all users and evaluates Bit Error Rate (BER) across several modulation schemes against brute-force (BF) ML and SIC. The simulation results show near-optimal detection performance, outperforming SIC and closely approaching exhaustive ML in the studied configurations. However, QA execution time is longer than both BF and SIC. These results set a baseline for future integrated quantum-classical systems within centralized RAN architectures.
\end{abstract}
